\begin{claim}{C1.2}{Proved}{A basis on which all first differentials are specified}
Choose a nonzero $g\in H^{0,1}_{\bar\partial}(X)$, a basis
$q_1,\ldots,q_h$ of a complement to $\C d_1g$ in
$H^{1,1}_{\bar\partial}(X)$, a basis $v_1,\ldots,v_h$ of
$H^{1,2}_{\bar\partial}(X)$, and a class $e$ complementary to the
image of $d_1:H^{1,2}\to H^{2,2}$.  With
$a=d_1g$, $c_j=d_1q_j$, $d_j=d_1v_j$, and $t=d_1e$, the first
differential is
\[
g\mapsto a\mapsto0,\quad q_j\mapsto c_j\mapsto0,
\qquad v_j\mapsto d_j\mapsto0,\quad e\mapsto t\mapsto0.
\]
These are bases in the corresponding bidegrees.  If $h=1$ and a
class $q$ restricts to the invariant fibre class $\eta$ described below,
we may take $q_1=q$, $v_1=gq$, and $e=q^2$, after scaling $g$ if
necessary.  Statements A1.3 and A3.2 prove these hypotheses for $X$,
so this geometric choice applies to the threefold in the question.
\end{claim}
\begin{proof}
The first assertions follow from the injectivity and surjectivity
proved in C1.1.  For the last assertion write
$\eta=u\wedge w+6\gamma\wedge\delta$ in the marked integral
cohomology of a regular fibre.  Its complex-oriented square has
integral $-12$.  Its product with the $(0,1)$ image of $\gamma$ is
nonzero: the explicit period calculation in C3.7 gives its nonzero
pairing with $d\zeta_2$.  Thus $gq$ and $q^2$ have nonzero fibre
restrictions.  The images under $d_1$ restrict to zero on the fibre,
because the fibre is a complex torus and its invariant representatives
are $d$-closed.  Therefore $q$ complements $d_1g$, $gq$ spans the
one-dimensional $H^{1,2}$, and $q^2$ complements $d_1H^{1,2}$.
The extension used in this last choice is the actual Leray lift in
Statement A3.2.  This argument uses that extension; it does not infer
it from a calculation on the smooth fibre alone.
\end{proof}
